%% LyX 2.4.4 created this file.  For more info, see https://www.lyx.org/.
%% Do not edit unless you really know what you are doing.
\documentclass[english]{article}
\usepackage[T1]{fontenc}
\usepackage[latin9]{inputenc}
\usepackage{units}
\usepackage{enumitem}
\usepackage{amsmath}
\usepackage{amsthm}
\usepackage{graphicx}
\usepackage{geometry}
\geometry{verbose,tmargin=1in,bmargin=1in,lmargin=1in,rmargin=1in}
\PassOptionsToPackage{normalem}{ulem}
\usepackage{ulem}

\makeatletter

%%%%%%%%%%%%%%%%%%%%%%%%%%%%%% LyX specific LaTeX commands.
\newcommand{\lyxmathsym}[1]{\ifmmode\begingroup\def\b@ld{bold}
  \text{\ifx\math@version\b@ld\bfseries\fi#1}\endgroup\else#1\fi}


%%%%%%%%%%%%%%%%%%%%%%%%%%%%%% Textclass specific LaTeX commands.
\numberwithin{equation}{section}
\numberwithin{figure}{section}
\newlength{\lyxlabelwidth}      % auxiliary length 

%%%%%%%%%%%%%%%%%%%%%%%%%%%%%% User specified LaTeX commands.
\usepackage{fancyhdr}
\usepackage{lastpage}
\pagestyle{fancy}
\usepackage{enumitem}
\usepackage{ifthen}
\everymath=\expandafter{\the\everymath\displaystyle}
%\DeclareMathSizes{10}{10}{10}{10}
\usepackage{multicol}

\usepackage{tikz}
\usepackage{tikz-3dplot}
\usetikzlibrary{calc,arrows}

\usetikzlibrary{patterns}
\usetikzlibrary{plotmarks}
\usepackage{pgfplots}
\pgfplotsset{compat=newest}
\usepgfplotslibrary{polar}

\usepackage{calc}
\usepackage[nomessages]{fp}

\usepackage{datetime}
% Added by lyx2lyx
\setlength{\parskip}{\smallskipamount}
\setlength{\parindent}{0pt}

\makeatother

\usepackage{babel}
\begin{document}
\renewcommand{\author}{Dr.\,James Ashe}

\newcommand{\classNum}{331}

\newcommand{\classSec}{A}

\newcommand{\examType}{Final Exam review}

\renewcommand{\date}{}

\newcommand{\name}{Final Exam time: \formatdate{3}{12}{2013}~ from 5--7pm}

%%First page's header%%%%%%%%%%%%%%%%%%%%%%
\fancypagestyle{firststyle} %%header settings for first pagr
{
	\fancyhead[L]{MTH \classNum{}(\classSec{}) \author{}\\
	\textbf{\examType{}} ~\date{}}
	%%%\fancyhead[C]{\textbf{\name}}
	\fancyhead[R]{\textbf{\name}}
	\setlength{\headsep}{14pt}
}
%%%%%%%%%%%%%%%%%%%%%%%%%%%%%%%%%%%%%%%%%%

%%Next page's header%%%%%%%%%%%%%%%%%%%%%%%
\setlist{nolistsep}
\lhead{MTH \classNum{}(\classSec{})} 
%\chead{\textbf{\name}}
\rhead{\textbf{\name}} 
\cfoot{\thepage{}~of~\pageref{LastPage}} 
\setlength{\headsep}{5pt} 
%%%%%%%%%%%%%%%%%%%%%%%%%%%%%%%%%%%%%%%%%%%

%%Various settings; see the preamble for other settings%%%%%
%\setenumerate[0]{leftmargin=12pt,itemindent=0pt,labelwidth=0pt}
%\setlist{topsep=.5pt}
\renewcommand{\labelenumi}{\textbf{\arabic{enumi}.}}
\renewcommand{\labelenumii}{\textbf{\alph{enumii}.}}
\thispagestyle{firststyle}
%%%%%%%%%%%%%%%%%%%%%%%%%%%%%%%%%%%%%%%%%%

\newcommand{\xmax}{0}
\newcommand{\xmin}{-\xmax}
\newcommand{\xstep}{0}
\newcommand{\ymax}{0}
\newcommand{\ymin}{-\ymax}
\newcommand{\ystep}{0} 
\newcommand{\dspace}{\vspace{1cm}}
\newcommand{\xa}{0}
\newcommand{\xb}{0}
\newcommand{\ya}{1}
\newcommand{\yb}{1}

\newcommand{\xspace}{\vspace{1cm}}

%%Various settings; see the preamble for other settings%%%%%
%%\setenumerate[0]{leftmargin=12pt,itemindent=0pt,labelwidth=0pt}
\renewcommand{\labelenumi}{\textbf{\arabic{enumi}.}}
\renewcommand{\labelenumii}{\textbf{\alph{enumii}.}}
\thispagestyle{firststyle}
%%%%%%%%%%%%%%%%%%%%%%%%%%%%%%%%%%%%%%%%%%

\subsubsection*{Reading and writing sigma notation}
\begin{enumerate}[resume]
\item \small Write the infinite series, $0.3+0.03+0.003+\cdots$, using
sigma notation (i.e. using $\Sigma$) . 
\item Write the first 3 terms of the sequence of partial sums of the series
${\displaystyle \sum_{k=1}^{\infty}\frac{3^{k}}{2}}$.
\end{enumerate}

\subsubsection*{Convergence/Divergence of series}

\small \emph{Evaluate the geometric series or state that it diverges. }
\begin{enumerate}[resume]
\item ${\displaystyle \sum_{k=0}^{\infty}\frac{2^{k}}{7^{k}}}$
\item ${\displaystyle \sum_{m=2}^{\infty}\left(\frac{1}{3}\right)^{m}}$
\end{enumerate}
\emph{Determine the radius of convergence of the following power series. }
\begin{enumerate}[resume]
\item ${\displaystyle \sum\frac{x^{k}}{2^{k}}}$
\end{enumerate}
\emph{Using the Squeeze Theorem, find the limit of the following sequence
or determine that the limit does not exist.}
\begin{enumerate}[resume]
\item ${\displaystyle \lim_{n\rightarrow\infty}\left\{ \frac{\sin^{2}n}{2^{n}}\right\} }$
\end{enumerate}
\emph{Use the Divergence Test to determine whether the following series
diverge, or state that the test is inconclusive.}
\begin{enumerate}[resume]
\item ${\displaystyle \sum_{k=1}^{\infty}\frac{k^{4}-1}{k^{4}+10}}$
\end{enumerate}
\emph{Use the Ratio Test to determine whether the following series
converges.}
\begin{enumerate}[resume]
\item ${\displaystyle \sum_{k=1}^{\infty}\frac{3k+3}{6(k+1)}}$
\end{enumerate}

\subsubsection*{Taylor polynomials}
\begin{enumerate}[resume]
\item \small Let $f(x)=2\cos(x)$. Find the Taylor polynomial of order
$n=3$ for $f(x)$ centered at $0$.

\begin{enumerate}[resume]
\item Use the Taylor polynomial from above to approximate $2\cos(0.1)$.
\item Compute the absolute error in the approximation assuming the exact
value is given by a calculator.
\end{enumerate}
\end{enumerate}

\subsubsection*{Parametric functions}

\small \emph{Determine $\nicefrac{dy}{dx}$ in terms of $t$ and
evaluate it at the given value of $t$. }
\begin{enumerate}[resume]
\item $x=2t$, $y=4t^{2}$; $t=1$.
\end{enumerate}
\emph{Give a set of parametric equations that describes the curve.}
\begin{enumerate}[resume]
\item A circle centered at $(0,0$) with radius 6, generated clockwise.
\end{enumerate}

\subsubsection*{Calculus in polar coordinates}

\small \emph{Locate the following polar coordinates on a graph. }
\begin{enumerate}[resume]
\item A:$\left(2,\frac{\pi}{4}\right)$, B:$\left(-2,\frac{\pi}{4}\right)$,
C:$\left(2,\frac{\pi}{4}+\pi\right)$, D:$\left(2,\frac{\pi}{4}+2\pi\right)$
\end{enumerate}
\emph{Express the following }\textbf{\emph{polar}}\emph{ coordinates
in }\textbf{\emph{Cartesian}}\emph{ coordinates.}
\begin{enumerate}[resume]
\item $\left(-4,\frac{2\pi}{3}\right)$
\end{enumerate}
\emph{Express the following }\textbf{\emph{Cartesian}}\emph{ coordinates
in }\textbf{\emph{polar}}\emph{ coordinates in at least }\textbf{\emph{two}}\emph{
different ways. }
\begin{enumerate}[resume]
\item $\left(-1,1\right)$
\end{enumerate}
\emph{Find the slope of the line tangent to the following polar curves
at the given point using calculus. }
\begin{enumerate}[resume]
\item $r=2\cos2\theta$; $\left(2,\frac{\pi}{2}\right)$.
\end{enumerate}

\subsubsection*{Vectors}

\small \emph{Let $\mathbf{u}=\left\langle 2,4,1\right\rangle $ and
$\mathbf{v}=\left\langle 1,-3,2\right\rangle $. Perform the indicated
operations.}
\begin{enumerate}[resume]
\item $\mathbf{u}-2\mathbf{v}$
\item $\mathbf{u\cdot}\mathbf{v}$
\item $\mathbf{u\times v}$
\end{enumerate}
\emph{Use the graph to draw the following vectors (provide labels).
Make certain to indicate direction.}
\begin{enumerate}[resume]
\item \begin{multicols}{2}

\begin{enumerate}[resume]
\item \textbf{$2\mathbf{u}$}
\item $\mathbf{u}+\mathbf{v}$ 
\item $\mathbf{u}-\mathbf{v}$ (you do not have to draw it as a position
vector)
\item $\mbox{proj}_{\mathbf{v}}\mathbf{u}$ 
\end{enumerate}
\columnbreak\includegraphics[width=2cm]{D:/home/jim/JamesAshe/JCSU/Calc_3_MTH_331/slides_Calc3/images/vectors_2D.svg}\end{multicols}
\end{enumerate}
\emph{Use a graph to perform the indicated operations (provide labels). }
\begin{enumerate}[resume]
\item Plot the point $\left(3,3,3\right)$.

\begin{enumerate}[resume]
\item Sketch the plane $\left(0,3,0\right)$.
\end{enumerate}
\end{enumerate}
\emph{Perform the indicated operation.}
\begin{enumerate}[resume]
\item Let $\mathbf{u}=5\mathbf{i+4}\mathbf{j}$ and $\mathbf{v=}3\mathbf{i}-\mathbf{j}$.
Find the angle between $\mathbf{u}$ and $\mathbf{v}$.
\item Find a vector perpendicular to $\mathbf{u}=\left\langle 2,0,0\right\rangle $
and $\mathbf{v}=\left\langle 0,-3,0\right\rangle $
\item A bullet is fired with a muzzle velocity of 1037 ft/sec from a gun
aimed at an angle of $8\textdegree$ above the horizontal. Find a
vector describing the vertical and horizontal components of the bullet's
velocity.
\end{enumerate}

\subsubsection*{Derivatives of multi-variable functions}

\small \emph{Find the indicated partial derivative of the following
functions.}
\begin{enumerate}[resume]
\item Find $g_{x}$ and $g_{y}$ for $g(x,y)=(x^{2}+2)e^{y}$.
\item Find the four second partial derivatives of $h(x,y)=x^{2}\cos2y$.
\end{enumerate}
\emph{Use the Chain Rule to find the following derivative. When feasible
express your answer in terms of the independent variable. }
\begin{enumerate}[resume]
\item $dz/dt$ where $z=x^{3}y-x$, $x=t^{2}$, and $y=2t^{-1}$.
\end{enumerate}

\subsubsection*{Maximum/Minimum of multi-variable functions}

\small \emph{Compute the gradient of the following function and evaluate
it at the given point P.}
\begin{enumerate}[resume]
\item $f(x,y)=x^{2}+y^{2}$; $P(-1,2)$
\end{enumerate}
\emph{Consider the following functions and points P. }

\textbf{a.}\emph{ Find the }\emph{\uline{unit}}\emph{ vectors that
give the direction of steepest ascent and steepest descent at P.}

\textbf{b.}\emph{ Find a vector that points in a direction of no
change in the function at P.}
\begin{enumerate}[resume]
\item $f(x,y)=x^{2}+y^{2}$; $P(-1,2)$
\end{enumerate}
\emph{Use the Second Derivative Test to determine (if possible) whether
each critical point of the given function is a local maximum, local
minimum, or saddle point.}
\begin{enumerate}[resume]
\item $f(x,y)=4+x^{3}+y^{3}-3xy$ with critical points $(0,0)$ and $(1,1).$
\end{enumerate}
\emph{Find all of the critical points (there are two) of the following
function.}
\begin{enumerate}[resume]
\item $f(x,y)=\frac{x^{3}}{3}-\frac{y^{3}}{3}+3xy$
\end{enumerate}

\subsubsection*{Lagrange Multipliers (bonus)}

\small \emph{Use Lagrange multipliers to find the maximum and minimum
values of $f$ (when they exist) subject to the constraint.}
\begin{enumerate}[resume]
\item $f(x,y)=x+2y$ subject to $x^{2}+y^{2}=4$
\item $f(x,y)=e^{2xy}$ subject to $x^{3}+y^{3}=16$
\end{enumerate}

\end{document}
